\section{Path moments, entropy and risk}
\label{sec:moments}
Throughout this section, first fix an embedding. All bounds will be uniform over embeddings, so they also hold for their mixture.

\begin{lemma}[Selected first moments and native fractional moments]
\label{lem:path-moments}
Let $S_\ell$ be the largest $\E R_j^\ell$ over selected hidden neurons, and put
\[
 P_\ell=\frac1N\sum_{j=1}^N\E(T_j^\ell)^{\alpha_{\ell+1}}
 \qquad(1\le\ell\le H).
\]
Under the comparison law,
\begin{align}
 S_\ell&\le C_\ell(d+1)(B+6)D^{\ell-1},\\
 P_\ell&\le C_\ell\left[(d+1)
       +\{(d+1)(B+6)D^{\ell-1}\}^{\alpha_{\ell+1}}\right].
 \label{eq:path-moments}
\end{align}
where $C_\ell$ may depend on depth, the tail indices and the activation, but not on $d,m,B,N,\epsilon$.
\end{lemma}
\begin{proof}
A selected first-layer row satisfies $R_j^1\le b+A(d+1)(B+1)$. For an unselected row, let $p=\alpha_2<\min\{1,\alpha_1\}$. Subadditivity and the finite $p$th coordinate moment give
\[
 \E(R_j^1)^p
 \le C\{1+(d+1)\E|W_{jk}^1|^p\}
 \le C\{1+(d+1)^{1-p/\alpha_1}\}\le C(d+1).
\]
This proves the initial case, including $P_1$.

For a later selected row, conditional on preceding layers, Lemma~\ref{lem:small-coordinate} gives
\begin{align*}
 \E\sum_{\mathrm{background}}|W_{jk}^\ell|R_k^{\ell-1}
 &\le \frac{C\epsilon^{1-\alpha_\ell}}N
          \sum_k(T_k^{\ell-1})^{\alpha_\ell}.
\end{align*}
Here and below conditional expectations are with respect to the current incoming row. Designated coefficients have absolute value at most $B+5$. Therefore
\begin{equation}
 S_\ell\le b+Am(B+5)S_{\ell-1}+CA P_{\ell-1},
 \label{eq:selected-first-moment}
\end{equation}
using $\epsilon\le1$.

At an unselected row, take $p=\alpha_{\ell+1}<\alpha_\ell<1$. Lemma~\ref{lem:absolute-sum} conditionally on preceding layers gives
\[
 \E[(T_j^\ell)^p\mid\text{preceding layers}]
 \le C\left[1+
   \left(\frac1N\sum_k(T_k^{\ell-1})^{\alpha_\ell}\right)^{p/\alpha_\ell}
 \right].
\]
Jensen's inequality applies because $p/\alpha_\ell<1$. At a selected row, $\E(T_j^\ell)^p\le1+S_\ell^p$. Consequently
\begin{equation}
 P_\ell\le C\{1+P_{\ell-1}^{p/\alpha_\ell}+S_\ell^p\}.
 \label{eq:fractional-moment}
\end{equation}
Induction first bounds $S_\ell$ by the displayed geometric majorant, because all powers in the fractional recurrence are below one. To retain the sharper bound on $P_\ell$, raise its preceding bound to $p/\alpha_\ell$: the teacher term becomes $\{(d+1)(B+6)D^{\ell-2}\}^{p}$, while the dimension term is at most $d+1$. Together with $S_\ell^p$, these are bounded by~\eqref{eq:path-moments}.
\end{proof}

The strict tail decrease has a precise use here. We request a moment of order $\alpha_{\ell+1}$ from a native layer of index $\alpha_\ell$, never a moment at its own tail index. Selected nodes have stronger moments because their incoming background coordinates are conditionally truncated.

\begin{lemma}[The comparison entropy]\label{lem:comparison-entropy}
For the comparison measure, uniformly over every finite $N\ge m$,
\begin{equation}
 \KL(\nu\|\mu_N)
 \le C\left[
 U_{m,B}\epsilon^{-a}
 +(d+1)M\log\frac{CM(d+1)(B+6)}\epsilon
 \right],
 \qquad a=\alpha_2.
 \label{eq:comparison-entropy}
\end{equation}
\end{lemma}
\begin{proof}
At a background coordinate entering layer $\ell$, its conditional KL is
$-\log\Pp\{|W|\le\epsilon/T_k^{\ell-1}\}$, bounded by
$C\epsilon^{-\alpha_\ell}(T_k^{\ell-1})^{\alpha_\ell}/N$.
Thus a selected row costs at most $C\epsilon^{-\alpha_\ell}P_{\ell-1}$ after expectation.

The number of selected rows entering layer $\ell\ge2$ is at most $m^{H-\ell+1}$, with one row at $\ell=H+1$. Since $\alpha_\ell\le a$ and $\epsilon\le1$, the total background cost is bounded by
\[
 C\epsilon^{-a}\sum_{\ell=2}^{H+1}m^{H-\ell+1}
 \left[(d+1)+\{(d+1)(B+6)D^{\ell-2}\}^{\alpha_\ell}\right].
\]
Use $D\le C_A m(B+6)$, $\alpha_\ell<1$, and $d+1\ge1$. The second summand in each term is bounded by
\[
 C(d+1)m^{H-\ell+1+\alpha_\ell(\ell-2)}
       (B+6)^{\alpha_\ell(\ell-1)}
 \le C(d+1)(m+1)^{H-1}(B+6)^{aH}.
\]
The first summands satisfy the same bound. Summing the fixed number of layers gives $C U_{m,B}\epsilon^{-a}$.

Use the relative-entropy chain rule for the feedforward conditional law. Lemma~\ref{lem:embedding-entropy} cancels the designated-edge $\log N$ terms in the large-width branch, while $\log N\le\log M$ suffices in the small-width branch. There are at most $M$ designated internal/readout edges and $M(d+1)$ first-layer coordinates. Their interval density costs give the logarithmic term in~\eqref{eq:comparison-entropy}. The moment bounds show that every conditional entropy expectation used in this chain rule is finite.
\end{proof}

\begin{lemma}[Pointwise raw perturbation moments]\label{lem:raw-perturbation}
Let $z_g$ denote the comparator's raw readout. Uniformly over the input cube,
\begin{equation}
 \E_\nu|z_w(x)-z_g(x)|^2
 \le V_{m,B}\epsilon^{2-a}.
 \label{eq:raw-rms}
\end{equation}
If $\phi'$ exists and is globally Lipschitz, the constant in $V_{m,B}$ may be enlarged so that also
\begin{equation}
 \left|\E_\nu[z_w(x)-z_g(x)]\right|
 \le V_{m,B}\epsilon^{2-a}.
 \label{eq:raw-bias}
\end{equation}
\end{lemma}
\begin{proof}
Teacher hidden values satisfy
\[
 \sup_x|h_j^{*,\ell}(x)|
 \le C(d+1)(B+6)D^{\ell-1}.
\]
Fix an input and let $v_\ell$ be the largest selected-neuron second moment of the error. In the first layer,
$v_1\le A^2(d+1)\epsilon^2/3$.

At a later selected node, write the preactivation error as
\[
 \Delta=\sum_{j=1}^m c_j(h_j-h_j^*)+Z.
\]
Conditionally on preceding layers, $Z$ is centered, and its variance is bounded by
\begin{equation}
 \E[Z^2\mid\text{preceding layers}]
 \le17\epsilon^2\sum_{\mathrm{designated}}|h_j|^2
 +\frac{C\epsilon^{2-\alpha_\ell}}N
       \sum_k(T_k^{\ell-1})^{\alpha_\ell}.
 \label{eq:adaptive-variance}
\end{equation}
For a background term, this follows because its conditional variance is at most
$C(\epsilon/T_k)^{2-\alpha_\ell}/N$, while $|h_k|^2\le T_k^2$.
No second moment of an unselected hidden feature is required.

Conditional centering removes the cross term with the preceding-layer error. Since
$|h_j|^2\le2|h_j-h_j^*|^2+2|h_j^*|^2$, Lipschitzness yields
\begin{align}
 v_\ell\le A^2\bigg[&(m^2B^2+34m\epsilon^2)v_{\ell-1}\\
 &+Cm(d+1)^2(B+6)^2D^{2\ell-4}\epsilon^2
 +C P_{\ell-1}\epsilon^{2-\alpha_\ell}\bigg].
 \label{eq:variance-recursion}
\end{align}
Induction, $\epsilon^2\le\epsilon^{2-a}$, $\alpha_\ell\le a$, and Lemma~\ref{lem:path-moments} give the following bound. In the active-noise term, $m(B+6)^2\le D^2$ absorbs the apparent extra coefficient factor; the first-layer error itself has no such factor:
\[
 v_\ell\le C_\ell(d+1)^2D^{2\ell-2}\epsilon^{2-a}.
\]
The readout obeys the same recursion with activation Lipschitz constant one. Its additional factor is bounded by a constant times $D^2$, proving~\eqref{eq:raw-rms}.

For the smooth case, let $J$ be a Lipschitz constant for $\phi'$, and let $b_\ell$ be the largest absolute selected-neuron mean error. Centering of the first-layer preactivation error gives
$b_1\le J(d+1)\epsilon^2/6$. Taylor's inequality at a later teacher preactivation gives
\[
 b_\ell\le AmB\,b_{\ell-1}+\frac J2\E\Delta^2.
\]
The already established variance bounds control the second term, yielding
$b_\ell\le C_\ell(d+1)^2D^{2\ell-2}\epsilon^{2-a}$.
The readout bias is at most $mB\,b_H$, because all its coefficient perturbations and adaptive backgrounds are conditionally centered. This proves~\eqref{eq:raw-bias}. Every estimate is uniform over embeddings and hence holds for their mixture.
\end{proof}

\begin{lemma}[A data-independent mean-square comparison]\label{lem:l2-comparison}
For each $g\in\calG_{m,B}$, $0<r\le1$ and finite $N\ge m$, there is a measure $\nu_{g,r,N}\ll\mu_N$, independent of the observations, satisfying
\begin{equation}
 \sup_x\E_\nu|f_w(x)-g(x)|^2\le r,\qquad
 \KL(\nu\|\mu_N)\le K_{m,B}r^{-q},
 \qquad q=\frac a{2-a}.
 \label{eq:l2-comparison}
\end{equation}
No smoothness assumption is needed.
\end{lemma}
\begin{proof}
A Lipschitz output link gives $\E|f_w-g|^2\le L_\psi^2V_{m,B}\epsilon^{2-a}$ by Lemma~\ref{lem:raw-perturbation}. Choose $\epsilon=(r/(L_\psi^2V_{m,B}))^{1/(2-a)}$, increasing constants so $\epsilon\le1$. Substitute into~\eqref{eq:comparison-entropy}. The power term is $C U_{m,B}V_{m,B}^{a/(2-a)}r^{-q}$, and $r^q\log(1/r)\le1/(eq)$ absorbs the logarithmic term. This is precisely the resource~\eqref{eq:resources}.
\end{proof}

\begin{lemma}[A data-independent risk comparison]\label{lem:comparison}
For each $g\in\calG_{m,B}$, every $0<r\le1$ and every finite $N\ge m$, there is $\nu_{g,r,N}\ll\mu_N$, independent of the observations, such that
\begin{equation}
 \E_\nu(f_w(x)-y)^2\le(g(x)-y)^2+r
 \quad(x\in[-1,1]^d,\ |y|\le1),\qquad
 \KL(\nu\|\mu_N)\le K_{m,B}r^{-q},
 \label{eq:risk-comparison}
\end{equation}
with exponent $q=2a/(2-a)$ in general, or $q=a/(2-a)$ if $\phi\in C^{1,1}$ and the output link has the sharper regularity specified in Section~2.
\end{lemma}
\begin{proof}
For a Lipschitz output link with constant $L_\psi$, the loss
$\ell_y(z)=(\psi(z)-y)^2$ is $4L_\psi$-Lipschitz for $|y|\le1$.
Thus~\eqref{eq:raw-rms} gives risk increase at most
$4L_\psi\sqrt{V_{m,B}}\,\epsilon^{(2-a)/2}$.
Choose $\epsilon$ to make this at most $r$ and use~\eqref{eq:comparison-entropy}. Its power term becomes
\[
 C U_{m,B}V_{m,B}^{a/(2-a)}r^{-2a/(2-a)}.
\]

For the sharper case, squared loss after the output link satisfies
\[
 \ell_y(z+u)\le\ell_y(z)+d_y(z)u+C_\psi u^2,
 \qquad |d_y(z)|\le D_\psi,
\]
with fixed constants. For clipping, one may use $D_\psi=4,C_\psi=1$, because $\ell_y(z)-z^2$ is concave. For a bounded smooth link with $|\psi'|\le L_\psi$ and $\operatorname{Lip}(\psi')\le J_\psi$, take
$D_\psi=4L_\psi$ and $C_\psi=L_\psi^2+2J_\psi$.
Equations~\eqref{eq:raw-rms}--\eqref{eq:raw-bias} therefore give risk increase at most $C V_{m,B}\epsilon^{2-a}$. The chosen cutoff now makes the entropy power term
\[
 C U_{m,B}V_{m,B}^{a/(2-a)}r^{-a/(2-a)}.
\]
In both cases the same resource~\eqref{eq:resources} suffices. The remaining logarithm of $1/r$ is absorbed using $r^q\log(1/r)\le1/(eq)$. All cutoff choices are at most one after increasing the fixed constants.
\end{proof}

The pointwise comparison concerns bounded reference labels. To handle unbounded observed labels, the next section applies it with $y=f_0(x)$, then concentrates the comparison-averaged relative loss.
